\begin{proof}
In the Laurent polynomial ring $K[x^{\pm1},y^{\pm1}]$, the prime elements
$x-1$ and $y-1$ have no common irreducible factor.  Since this ring is a
UFD,
\[
 ((x-1)(y-1))=(x-1)\cap(y-1).
\]
Hence $X$ is the reduced union $Y_h\cup Y_v$.  Both components are
isomorphic to $\mathbf G_m^{\an}$, and their intersection is the rigid
point $p=(1,1)$.  Algebraic analytification over a nontrivially valued
field gives a good strict boundaryless space; here every irreducible
component has dimension one, so $X$ is pure one-dimensional.

The functions $x$ and $y$ are invertible on $X$.  Thus the coordinates of
$\tau$ are integral harmonic locally.  On $Y_h$, the map is
\[
 -\log|x|\longmapsto(-\log|x|,0),
\]
and on $Y_v$ it is
\[
 -\log|y|\longmapsto(0,-\log|y|).
\]
The moment skeleton of $\mathbf G_m^{\an}$ is the real line of Gauss
points, and each displayed lattice map sends $1$ to the corresponding
primitive basis vector.  The generic analytic degree and the generic
multiplicity of each reduced component are one.  It follows from the
tropical multiplicity formula, componentwise additivity, and the two
primitive lattice indices that the image is the cross and that its
fundamental weights are all one.

The map $\tau$ is proper over compact subsets: on either component the
preimage of a bounded interval is a closed annulus, and the preimage of a
compact subset of $\mathbf R^2$ is closed in a finite union of such compact
annuli.  If $\rho>1$ belongs to $|K^\times|$, the displayed inequalities
therefore define a compact strict analytic domain $X_\rho$, and its image is
\[
 [-\log\rho,\log\rho]\times\{0\}
 \ \cup
 \{0\}\times[-\log\rho,\log\rho].
\]
Because the coordinate functions are logarithms of analytic units, the
restriction is a harmonic tropicalization in the compact-domain sense.
The full noncompact map is an ordinary tropicalization that is locally
harmonic.

For $Z=2[Y_h]+[Y_v]$, additivity of the tropicalization of analytic cycles
and the preceding primitive index calculation give weights
$(2,2,1,1)$ in the directions $(e_1,-e_1,e_2,-e_2)$.  The balancing sum is
\[
 2e_1+2(-e_1)+e_2+(-e_2)=0.
\]
The integration current of an effective analytic cycle is strongly
positive and of order zero, and integration currents of analytic cycles on
a boundaryless space are closed.  Thus
$T=2\delta_{Y_h}+\delta_{Y_v}$ is positive and closed.  The analytic
tropical projection formula identifies its chartwise pushforward with the
balanced skeletal current having those weights.

The functions $u$ and $v$ are continuous and convex on $\mathbf R^2$.
Their pullbacks are therefore continuous psh functions.  More explicitly,
for $\varepsilon_j\downarrow0$ set
\[
 u_j(s,t)=2\sqrt{s^2+\varepsilon_j},\qquad
 v_j(s,t)=\sqrt{s^2+\varepsilon_j}
          +\sqrt{t^2+\varepsilon_j}.
\]
Since
\[
 \frac{\mathrm d^2}{\mathrm dr^2}\sqrt{r^2+\varepsilon_j}
 =\frac{\varepsilon_j}{(r^2+\varepsilon_j)^{3/2}}\geq0,
\]
these are smooth convex functions.  They decrease locally uniformly to
$u,v$, and their pullbacks are smooth psh functions decreasing locally
uniformly to $U,V$.

Let $q_h$ be the unit Gauss point in the $x$-coordinate on $Y_h$, and let
$q_v$ be the unit Gauss point in the $y$-coordinate on $Y_v$.  Both map to
the origin.  They are distinct because
\[
 \begin{array}{c|cc}
       &|x-1|&|y-1|\\ \hline
  q_h  &1&0\\
  q_v  &0&1.
 \end{array}
\]
On the horizontal component put $s=-\log|x|$ and on the vertical component
put $t=-\log|y|$.  On the skeleton of each copy of $\mathbf G_m^{\an}$,
the function $|s|$, respectively $|t|$, has outgoing slope one in both
directions at the unit Gauss point and is affine elsewhere.  The graph
Laplacian formula therefore gives
\[
 \dprimeddprime[|s|]=2\delta_{q_h},\qquad
 \dprimeddprime[|t|]=2\delta_{q_v}.
\]
There is no atom at the rigid intersection point $p$: on the neighbourhood
$|x-1|<1$, $|y-1|<1$ of $p$, the ultrametric inequality gives
$|x|=|y|=1$, so all the functions under consideration are locally
constant.  There is no term at infinity because currents are tested against
compactly supported forms.

The normalization is the disjoint union $Y_h\sqcup Y_v$, and its finite
map to $X$ is proper.  Integration against the fundamental cycle is
componentwise, and differentiation commutes with this proper pushforward.
The restrictions are
\[
 \begin{array}{c|cc}
       &Y_h&Y_v\\ \hline
 U     &2|s|&0\\
 V     &|s|&|t|\\
 H=U-V &|s|&-|t|.
 \end{array}
\]
The preceding Laplacian calculation now yields
\[
 \dprimeddprime[U]=4\delta_{q_h},\qquad
 \dprimeddprime[V]=2\delta_{q_h}+2\delta_{q_v}.
\]
Since $T$ has order zero, multiplication by the continuous function $H$ is
defined.  Componentwise,
\[
 HT=2|s|\,\delta_{Y_h}-|t|\,\delta_{Y_v},
\]
and hence
\[
 \dprimeddprime(HT)=4\delta_{q_h}-2\delta_{q_v}.
\]

Finally, the defining pairing for tropical pushforward gives
\[
 \tau_*\delta_{q_h}=\tau_*\delta_{q_v}=\delta_0.
\]
Pushing forward the three formulas above gives
\[
 \tau_*\dprimeddprime[U]=\tau_*\dprimeddprime[V]=4\delta_0,
 \qquad
 \tau_*\dprimeddprime(HT)=2\delta_0.
\]
The first two analytic currents are unequal because $q_h\ne q_v$.  This
proves both the realization claim and its final scope restriction.
\end{proof}

\begin{quote}
\textbf{Verification inputs.}
The following previously established facts are the only imported inputs.
Algebraic analytifications over a nontrivially valued field are good,
strict, and boundaryless.  Logarithms of absolute values of analytic units
are integral harmonic.  The moment skeleton of
$\mathbf G_m^{\an}$ is its real line of Gauss points.  Tropical
multiplicities obey the primitive-lattice-index formula and are additive in
analytic cycles.  Continuous convex functions on a tropical chart pull
back to continuous psh functions, and smooth convex functions pull back to
smooth psh functions.  Analytic cycle currents are positive, order zero,
and closed on a boundaryless space.  Proper pushforward commutes with
$\mathrm d'$ and $\mathrm d''$.  On a weighted metric graph, the atom of
$\mathrm d'\mathrm d''[f]$ is the weighted sum of outgoing slopes.  These
facts are stated and proved or source-imported in Chapters 12--18 of the
manuscript.  The verifier is asked to audit the deductions, hypotheses,
normalizations, and the distinction between analytic equality and equality
only after tropical pushforward.
\end{quote}
